\section*{Limitations}
Although \modelname shows promising results for generating attractive but faithful headlines, there are still
some limitations: (1)~\modelname is a monolingual model that only
supports Chinese. It requires three pre-trained scorers. Also, as the FR labels are
specifically difficult to obtain, it is not easy to implement in
other languages. (2)~Running the whole framework with a batch size of 16
takes around 22\,GB GPU memory, mostly because we must load all pre-trained
models into the GPU. This can be alleviated by using a distilled pre-trained model. (3)~On average, \modelname generates more faithful
headlines than other baselines, but it still occasionally produces false
information or unwanted results. This work is only for academic purposes and is
not ready for production.

\section*{Ethics Statement}
Given that our dataset is in Chinese and requires a profound understanding of forward referencing for annotation and evaluation, we carefully selected annotators from our lab who specialize in NLP-related research and possess knowledge in linguistics. To ensure fairness, we provided all annotators with a payment of \$6.66 per hour, which is 10\% higher than the minimum hourly wage requirement in Taiwan.

During the data annotation process, we introduced the concept of forward referencing to the annotators, along with relevant examples. Only annotators who achieved an accuracy rate of over 80\% were eligible to perform the actual annotation task. It's important to note that we solely asked annotators to label the "type of forward reference," which is well-defined, and not to assess the accuracy or truthfulness of the news articles. With five annotators who successfully passed the pretest, combined with the relatively objective nature of labeling forward reference types, we believe any potential bias during the data annotation process is minimal.

For the evaluation phase, an additional five annotators were tasked with determining the preferable title among two options, based on attractiveness, faithfulness, and fluency. These annotators are different from those who labeled the data to ensure a blind test. Although this task involves a greater level of subjectivity, we provided average statistics based on the assessments of the five annotators. Additionally, we maintained a blind test by recruiting separate evaluators and randomly shuffling the order of the two titles for each trial. This evaluation protocol aligns with standard practices employed in the research community, and we believe it effectively minimizes potential biases.

It is also important to note that we are not really \emph{learning to mimic fake news},
by taking fake news headlines as the ground truth reference. Instead, we seek to
learn the \emph{writing techniques} that are often used in fake news to attract
readers. As we are aware of the risk of producing misinformation, we want to
again highlight the importance of the faithfulness check. \modelname was 
designed only to assist journalists as a reference to write 
faithful headlines that users prefer for verified news. Even if we propose using a faithfulness
scorer to increase fidelity, its nature, similar to attractive headline
generation systems, still exhibits the risk that \modelname could be used by
malicious users to generate sensational headlines for fake news. 
Additionally, \modelname may misjudge offensive or unethical headlines to be
a headline that users would prefer. 
%For such unwanted purpose, though \modelname isn't designed for fake news, they can generate more user-desired fake news headlines by removing the faithfulness scorer during training.
Our goal is 
% to leverage
  to fight fire with fire by leveraging   % AMH: check
fake news as learning material to
fight against misinformation, by encouraging users to read verified news. We call on users not to abuse \modelname to produce false information.
%We're calling for the design of an ethic-aware headline generation system, for example, an open-source fact-checking system that can stores pairs of fake news and verified news, as a useful database to help improve fact-checking.
\section*{Acknowledgement}
This work is supported by the National Science and Technology Council of Taiwan under grants 111-2221-E-001-021 and 111-2634-F-002-022.